\chapter{Dominance and Loops}
\label{ch:dominance-loops}

Knowing the shape of the graph is not the same as understanding it.
An optimizer must answer sharper questions: does every path to a block pass through the loop test, and which blocks belong to the loop itself.
This chapter defines \term{dominance}, computes it with a simple fixed-point iteration, and uses it to identify loops exactly.
Every loop optimization in the rest of the book rests on the two definitions given here.

\section{Who must run first}
\label{sec:ch08-must-run-first}

Consider a test block that guards a loop body.
Intuitively the test runs before each body iteration on every possible execution.
Dominance makes that intuition precise: a block $A$ \term{dominates} a block $B$ when every path from the entry to $B$ passes through $A$.
We write $A \in \dom(B)$ for the set of dominators of $B$.
Every block dominates itself, since the trivial path from the entry that ends at $B$ trivially passes through $B$ only when reached, and the entry dominates everything because every path starts there.
The interesting cases sit between: in a loop, the head dominates the body, but the body never dominates the head.
Strict dominance excludes the block itself, and the \term{immediate dominator} $\idom(B)$ is the closest strict dominator of $B$.
Linking every block except the entry to its immediate dominator forms the \term{dominator tree}, which summarizes all must-pass-through facts in $N - 1$ edges for $N$ blocks.
Figure~\ref{fig:ch08-cfg-domtree} shows a five-block loop graph beside its dominator tree.
\begin{figure}
\centering
\begin{tikzpicture}[node distance=1.5cm]
\node[basic block] (e) {B1 entry};
\node[basic block, below of=e] (h) {B2 head};
\node[basic block, below left of=h, xshift=-0.8cm] (a) {B3 left};
\node[basic block, below right of=h, xshift=0.8cm] (b) {B4 right};
\node[basic block, below of=h, yshift=-2.2cm] (x) {B5 exit};
\draw[cfg edge] (e) -- (h);
\draw[cfg edge] (h) -- (a);
\draw[cfg edge] (h) -- (b);
\draw[cfg edge] (a) -- (x);
\draw[cfg edge] (b) -- (x);
\draw[back edge] (x) -- (h);
\end{tikzpicture}
\caption{A five-block loop graph whose head dominates both arms while neither arm dominates the other.}
\label{fig:ch08-cfg-domtree}
\end{figure}
The entry reaches the head, the head branches to two arms, both arms join at the exit block, and the exit returns to the head along the dashed edge.
The head dominates both arms because every path from the entry to either arm passes through it.
Neither arm dominates the other because each is reachable without the other.
The exit block is dominated by the head and the entry but by neither arm.
Table~\ref{tab:ch08-domsets} records the full dominator sets that Section~\ref{sec:ch08-worked-example} computes by hand.
\begin{table}
\caption{Dominator sets for the five-block graph of Figure~\ref{fig:ch08-cfg-domtree}.}
\label{tab:ch08-domsets}
\centering
\begin{tabular}{ll}
\toprule
Block & Dominators \\
\midrule
B1 entry & B1 \\
B2 head & B1, B2 \\
B3 left & B1, B2, B3 \\
B4 right & B1, B2, B4 \\
B5 exit & B1, B2, B5 \\
\bottomrule
\end{tabular}
\end{table}
Each row contains the block itself plus exactly the blocks that every entry path must cross.
The head appears in every row except the entry because all flow passes through it.
The arms appear only in their own rows, which is the set-theoretic way of saying they are optional relative to each other.
\section{Computing dominators}
\label{sec:ch08-computing-dominators}

Dominators are the greatest solution of a simple system of set equations.
The entry is dominated only by itself, and every other block is dominated by itself plus the blocks that dominate all of its predecessors.
Intuitively, a block $A$ survives in $\dom(B)$ only when $A$ equals $B$ or when every predecessor already contains $A$.
Algorithm~\ref{alg:ch08-dom} starts with an overestimate and removes blocks until nothing changes.
\begin{algorithm}
\caption{Iterative computation of dominator sets.}
\label{alg:ch08-dom}
\begin{pseudocode}
\Require control-flow graph with entry $E$ and reachable block set $V$
\Ensure $\dom(B)$ for every block $B$
\Function{Dominators}{$E, V$}
\State set $\dom(E)$ to $\{E\}$ and every other $\dom(B)$ to $V$
\State $changed \gets true$
\While{$changed$}
\State $changed \gets false$
\ForAll blocks $B$ other than $E$
\State $new \gets \{B\} \cup \bigcap_{P\in\Pred(B)}\dom(P)$
\If{$new \neq \dom(B)$} \State $\dom(B)\gets new$; $changed\gets true$ \EndIf
\EndFor
\EndWhile
\State \Return all $\dom(B)$ sets
\EndFunction
\end{pseudocode}
\end{algorithm}
The algorithm assumes every block is reachable from the entry along at least one path.
Exclude unreachable blocks and their edges first; an unreachable cycle can have nonempty predecessor sets and still have no path from entry.
It guarantees the greatest fixed point of the equations, which coincides with the path-based definition of dominance.
Starting from the full set $V$ everywhere means the first updates only remove blocks, never add them.
Removal is monotone: once a block leaves a set it never returns, so the loop must terminate.
Each iteration processes every edge once to intersect predecessor sets, exactly like the propagation pattern of the previous chapter.
Production compilers often replace this iteration with the Lengauer-Tarjan method, which is faster on large graphs.
% TODO(cite): add the primary reference for the Lengauer-Tarjan dominator algorithm when the bibliography grows.
The iterative version stays because it fits on one page, obviously matches the definition, and is fast enough for the graphs in this book.

\section{Natural loops}
\label{sec:ch08-natural-loops}

Dominance turns the informal idea of a loop into a checkable definition.
A \term{back edge} is an edge from a block $T$ to a block $H$ where $H$ dominates $T$.
The target $H$ is the \term{loop header} and the source $T$ is the \term{loop tail}.
The \term{natural loop} of that edge contains $H$, $T$, and every block that can reach $T$ without passing through $H$.
It is found by walking backward from $T$ along predecessors and stopping at $H$.
In Figure~\ref{fig:ch08-cfg-domtree} the dashed edge runs from $B5$ to $B2$, and Table~\ref{tab:ch08-domsets} shows $B2$ in $\dom(B5)$.
The edge therefore qualifies, with header $B2$ and tail $B5$.
Walking backward from $B5$ reaches $B3$ and $B4$ through the join, then $B2$ itself, so the natural loop is $\{B2, B3, B4, B5\}$.
That set matches the visual loop: the head, both arms, and the exit block that jumps back.
Some graphs defeat this definition.
An \term{irreducible} graph has a loop with two entries, so no single header dominates the whole loop.
No back edge qualifies and the loop has no natural loop under this definition.
Such graphs arise from unstructured jumps and from some generated code.
Compilers either restructure the graph first or analyze irreducible regions with more general methods.
Mini programs built from structured statements always produce reducible graphs, so every loop the Mini front end creates has a natural loop.

\section{Worked example}
\label{sec:ch08-worked-example}

We now run Algorithm~\ref{alg:ch08-dom} on Figure~\ref{fig:ch08-cfg-domtree} by hand.
The predecessor sets are $\Pred(B1) = \emptyset$, $\Pred(B2) = \{B1, B5\}$, and $\Pred(B3) = \{B2\}$.
In addition $\Pred(B4) = \{B2\}$ and $\Pred(B5) = \{B3, B4\}$.
Initialization sets $\dom(B1) = \{B1\}$ and every other set to all five blocks $\{B1, B2, B3, B4, B5\}$.
Round one updates in order: $\dom(B2)$ becomes $\{B2\}$ union the intersection of $\{B1\}$ and all five blocks, which is $\{B1, B2\}$.
Next $\dom(B3)$ becomes $\{B3\}$ union $\dom(B2)$, which is $\{B1, B2, B3\}$.
Next $\dom(B4)$ becomes $\{B4\}$ union $\dom(B2)$, which is $\{B1, B2, B4\}$.
Next $\dom(B5)$ becomes $\{B5\}$ union the intersection of $\dom(B3)$ and $\dom(B4)$, which share $\{B1, B2\}$, giving $\{B1, B2, B5\}$.
Round two recomputes $\dom(B2)$ as $\{B2\}$ union the intersection of $\dom(B1)$ and $\dom(B5)$.
Since $\dom(B1) = \{B1\}$ and $\dom(B5) = \{B1, B2, B5\}$, the result is $\{B1, B2\}$, unchanged.
The remaining blocks recompute identically: $B3$ stays $\{B1, B2, B3\}$, $B4$ stays $\{B1, B2, B4\}$, and $B5$ stays $\{B1, B2, B5\}$.
No set changed in round two, so the iteration stops after $2$ rounds with the sets of Table~\ref{tab:ch08-domsets}.
The immediate dominators read off directly: $\idom(B2) = B1$, $\idom(B3) = B2$, $\idom(B4) = B2$, and $\idom(B5) = B2$.
The dominator tree is therefore $B1$ above $B2$ above $B3$, $B4$, and $B5$ as three children.
The back-edge check then finds exactly one natural loop $\{B2, B3, B4, B5\}$, as claimed in Section~\ref{sec:ch08-natural-loops}.
The hand computation uses $5$ blocks, $6$ edges, $2$ rounds, and $8$ set intersections in total.

\section{Implementation notes}
\label{sec:ch08-implementation}

Dominator sets are usually bit vectors with one bit per block.
Intersection becomes a loop of machine-word AND operations, and equality becomes a word comparison.
Numbering blocks in reverse postorder makes predecessors settle before their successors on most graphs, so convergence takes very few rounds.
Two or three rounds suffice for structured code, and the book graphs converge in at most as many rounds as the deepest nesting.
The running time of the naive iteration is $O(R \cdot (N + E) \cdot W)$, where $R$ is the number of rounds, $N$ is the block count, $E$ is the edge count, and $W$ is the cost of one set intersection.
With bit vectors $W$ is $N$ divided by the machine word size, and $R$ is bounded by $N$ in the worst case but tiny in practice.
Memory use is $O(N^2)$ bits for one set per block, which is one bit vector of length $N$ for each of the $N$ blocks.
Immediate dominators are derived by comparing set sizes: $\idom(B)$ is the strict dominator with the largest set.
Natural loops are collected by the backward walk described in Section~\ref{sec:ch08-natural-loops}, visiting each loop edge once.
A common bug is running the iteration before deleting unreachable blocks, which lets dead predecessors shrink live dominator sets incorrectly.
Another is forgetting that the entry needs no predecessors: seeding $\dom(E)$ to $\{E\}$ and never updating it keeps the base case stable.
Later chapters consume these results directly: loop-invariant motion needs headers, construction of static single assignment needs dominance frontiers, and control dependence needs the postdominator mirror image.
\begin{ccwarning}[Dominance needs reachability]
The equations assume every block is reachable from the entry.
Delete unreachable blocks before iterating, or their phantom predecessors will corrupt the sets of live code.
\end{ccwarning}

\section{Exercises}
\label{sec:ch08-exercises}

\begin{enumerate}
\item \label{ex:ch08-recompute} Recompute Table~\ref{tab:ch08-domsets} with the update order $B5, B4, B3, B2$ instead of the order used in Section~\ref{sec:ch08-worked-example}. Give the sets after each round, count the intersections, and confirm the same fixed point in the same number of rounds.
\item \label{ex:ch08-find-loops} Add an edge from $B4$ to $B3$ in Figure~\ref{fig:ch08-cfg-domtree} and rerun Algorithm~\ref{alg:ch08-dom}. List the back edges, the header and tail of each, and the natural loop of each. State which loops are nested and which share their header.
\item \label{ex:ch08-irreducible} Draw a three-block graph with entry $B1$ branching to $B2$ and $B3$ plus edges $B2$ to $B3$ and $B3$ to $B2$. Compute all dominator sets, show that neither cycle edge qualifies as a back edge, and explain why the cycle has no natural loop.
\end{enumerate}

\section{Further reading}
\label{sec:ch08-further-reading}

Appel uses dominators and loops as the scaffolding for dataflow analysis and code motion, which shows why this chapter precedes all optimization chapters, see~\cite{appel-1998-modern}.
Morgan covers loop handling as a practical backend concern, where headers and back edges decide what the optimizer may safely move, see~\cite{morgan-1998-optimizing}.
Safonov treats such structural analyses as checkable claims about the program rather than folklore about its shape, see~\cite{safonov-2010-trustworthy}.
The C language definition specifies the loop statements whose lowering produces the back edges analyzed here, see~\cite{wg14-2011-c11}.
